\renewcommand{\Anno}{2008}

\renewcommand{\NumeroProblema}{1}

\renewcommand{\Testo}{
\begin{enumerate}
    \item Determinare quanti sono i numeri interi positivi che, nella usuale rappresentazione decimale, hanno \emph{esattamente} $k$ cifre, tutte diverse da $7$.

    \item Sia $S_k$ la somma dei reciproci di tutti i numeri interi positivi che, nella usuale rappresentazione decimale, hanno \emph{esattamente} $k$ cifre, tutte diverse da $7$.
    
    Dimostrare che $S_{k}\leq 8$ per ogni valore di $k$.

    \item Sia $T_k$ la somma dei reciproci di tutti i numeri interi positivi che, nella usuale rappresentazione decimale, hanno \emph{al più} $k$ cifre, tutte diverse da $7$.
    
    Determinare se esiste una costante $C$ tale che $T_k\le C$ per qualunque valore di $k$.
\end{enumerate}
}

\renewcommand{\Risposta}{
\begin{enumerate}
	\item[(a)]  I numeri sono $8\cdot 9^{k-1}$.

	\item[(b)]  \ldots

	\item[(c)]  Esiste e la più piccola costante è $C=80$.
\end{enumerate}
}

\renewcommand{\Aiutino}{
\begin{enumerate}
	\renewcommand{\labelenumi}{(\alph{enumi})} 
	\item[(a)]  In quanti modi possiamo scegliere la prima cifra (quella più a sinistra)? In quanti modi possiamo scegliere le restanti cifre?

	\item[(b)]  Da quanti addendi è composta la somma che costituisce $S_{k}$? Qual è il più grande di questi addendi? Pertanto $S_{k}\leq\ldots$

	\item[(c)]  Come sono legati i $T_{k}$ agli $S_{k}$? Utilizzando la maggiorazione trovata al punto precedente ci si riduce alla somma dei termini di una progressione geometrica.
\end{enumerate}
}

\renewcommand{\SoluzioneUfficiale}{\TitoloSoluzione \DomandaA
La prima cifra (quella più a sinistra) può essere scelta in 8 modi (deve essere diversa da $0$ e da $7$), mentre le altre possono essere scelte in 9 modi, in quanto devono essere diverse solo da $7$. Trattandosi di scelte indipendenti abbiamo pertanto $8\cdot 9^{k-1}$ tali numeri.

\DomandaB I numeri con $k$ cifre sono tutti maggiori od uguali a $10^{k-1}$,
quindi i loro reciproci sono minori od uguali a $10^{1-k}$.  Per quanto dimostrato al punto precedente $S_{k}$ è la somma di $8\cdot 9^{k-1}$ addendi, tutti minori od uguali a $10^{1-k}$. Ne segue che
$$S_{k}\leq\frac{8\cdot 9^{k-1}}{10^{k-1}}=
8\left(\frac{9}{10}\right)^{k-1}.$$

In particolare
$$S_{k}\leq
8\left(\frac{9}{10}\right)^{k-1}\leq 8,$$
con uguaglianza se e solo se $k=1$. 

\DomandaC Mostriamo che $T_{k}\leq 80$ per ogni valore di $k$.
Per quanto trovato al punto precedente si ha infatti che
$$T_{k}=S_{1}+S_{2}+\ldots+S_{k}\leq
8\cdot\left[1+\frac{9}{10}+\left(\frac{9}{10}\right)^{2}+
\ldots+\left(\frac{9}{10}\right)^{k-1}\right].$$

Per calcolare il termine tra parentesi quadre possiamo ricorrere al caso particolare $a=9/10$ della nota formula per la somma dei termini di una progressione geometrica
$$1+a+a^{2}+\ldots+a^{k-1}=\frac{1-a^{k}}{1-a},$$
valida per ogni $a\neq 1$. Si ha pertanto che
$$T_{k}\leq
8\cdot\frac{1-\displaystyle{\left(\frac{9}{10}\right)^{k}}}
{1-\displaystyle{\frac{9}{10}}}
\leq 8\cdot\frac{1}{1-\displaystyle{\frac{9}{10}}}= 80,$$
il che completa la dimostrazione.
}

\renewcommand{\SoluzioniAlternative}{}

\renewcommand{\Osservazioni}{
}

\renewcommand{\Commenti}{
}